\DocumentMetadata{
  pdfversion=2.0,
  pdfstandard=ua-2,
  lang=en-US,
  tagging=on,
  tagging-setup={math/setup=mathml-SE,tabsorder=structure}
}
\documentclass[11pt,letterpaper]{article}
\usepackage[margin=0.68in,includefoot]{geometry}
\usepackage{newcomputermodern}
\usepackage{amsmath,amscd,mathtools}
\usepackage{tikz,adjustbox}
\providecommand{\square}{\mdlgwhtsquare}
\usepackage[hidelinks]{hyperref}
\hypersetup{
  pdftitle={Algebra PhD Exam, May 2004},
  pdfauthor={University of Florida Department of Mathematics},
  pdfsubject={Graduate mathematics examination},
  pdfdisplaydoctitle=true,
  bookmarksopen=true
}
\setcounter{secnumdepth}{0}
\setlength{\parindent}{0pt}
\setlength{\parskip}{0.28em}
\setlength{\emergencystretch}{3em}
\setlength{\leftmargini}{2.15em}
\setlength{\leftmarginii}{2.65em}
\setlength{\leftmarginiii}{3em}
\pagestyle{empty}
\raggedbottom
\allowdisplaybreaks
\NewTaggingSocketPlug{block/list/label}{examproblem}{%
  \tagstructbegin{tag=Lbl}\tagstructend%
  \tagstructbegin{tag=\LBody}%
  \tagstructbegin{tag=\ProblemHeadingTag,title-o={\ProblemHeadingText}}%
  \tagmcbegin{tag=\ProblemHeadingTag,actualtext={\ProblemHeadingText}}%
  #2%
  \tagmcend%
  \tagstructend%
}
\newenvironment{examproblems}[2]{%
  \def\ProblemHeadingTag{#2}%
  \AssignTaggingSocketPlug{block/list/label}{examproblem}%
  \begin{enumerate}%
  \setcounter{enumi}{#1}%
  \renewcommand{\labelenumi}{\textbf{\theenumi.}}%
  \setlength{\topsep}{0.35em}%
  \setlength{\partopsep}{0pt}%
  \setlength{\itemsep}{0.48em}%
  \setlength{\parsep}{0.18em}%
}{\end{enumerate}}
\newenvironment{examsubparts}{%
  \AssignTaggingSocketPlug{block/list/label}{default}%
  \begin{enumerate}%
  \setlength{\topsep}{0.18em}%
  \setlength{\partopsep}{0pt}%
  \setlength{\itemsep}{0.16em}%
  \setlength{\parsep}{0.1em}%
}{\end{enumerate}}
\newenvironment{examinstructions}{%
  \par\begingroup\itshape
}{%
  \par\endgroup\vspace{0.18em}
}
\makeatletter
\renewcommand\subsection{\@startsection{subsection}{2}{0pt}%
  {-1.25ex plus -.25ex minus -.1ex}%
  {0.55ex plus .1ex}%
  {\normalfont\large\bfseries}}
\renewcommand{\@maketitle}{%
  \newpage\null\vskip -1em%
  {\footnotesize\bfseries
    \href{https://www.ufl.edu/}{University of Florida}, \href{https://math.ufl.edu/}{Department of Mathematics}\par}%
  \vskip -0.5em%
  \tagstructbegin{tag=H1,title={Algebra PhD Exam, May 2004}}%
  \tagpdfparaOff%
  \tagmcbegin{tag=H1}%
    {\LARGE\bfseries\href{https://gma.math.ufl.edu/exams/algebra-phd/}{Algebra PhD Exam}, May 2004\par}%
  \tagmcend%
  \tagpdfparaOn%
  \tagstructend%
  \vskip 0.25em%
  \tagmcbegin{artifact}%
    \hrule height 1pt%
  \tagmcend%
  \vskip 1em%
}
\makeatother
\title{Algebra PhD Exam, May 2004}
\author{}
\date{}
\begin{document}
\maketitle
\thispagestyle{empty}
\begin{examinstructions}
Answer seven problems. Write your answers clearly in complete English sentences. You may quote results (within reason) as long as you state them clearly.
\end{examinstructions}
\begin{examproblems}{0}{H2}
\def\ProblemHeadingText{Problem 1.}
\item
Let~\(F\) be a field and let~\(E\) be a finite extension of~\(F\). Say that \(E/F\) is a simple extension if there is \(\alpha\in E\) such that \(E=F(\alpha)\).
\begin{examsubparts}
\item[\textnormal{(a)}]
Prove that \(E/F\) is a simple extension if and only if there are only finitely many fields~\(K\) such that \(E\supset K\supset F\).
\item[\textnormal{(b)}]
Prove that if~\(F\) has characteristic~\(0\) then \(E/F\) is a simple extension.
\end{examsubparts}
\def\ProblemHeadingText{Problem 2.}
\item
Let~\(p\) and~\(q\) be primes and let \(E\subset\mathbb C\) be the splitting field for \(g(X)=X^p-q\) over \(\mathbb Q\).
\begin{examsubparts}
\item[\textnormal{(a)}]
Give generators for~\(E\) as an extension of \(\mathbb Q\).
\item[\textnormal{(b)}]
Determine the degree of the extension \(E/\mathbb Q\).
\item[\textnormal{(c)}]
Give generators and relations for \(\operatorname{Gal}(E/\mathbb Q)\).
\item[\textnormal{(d)}]
Describe how the generators of \(\operatorname{Gal}(E/\mathbb Q)\) act on the generators of~\(E\).
\end{examsubparts}
\def\ProblemHeadingText{Problem 3.}
\item
Let~\(G\) be a finite nilpotent group, i.e., a group whose ascending central series has the form
\[
1=Z_0(G)\subsetneq Z_1(G)\subsetneq\cdots\subsetneq Z_{n-1}(G)\subsetneq Z_n(G)=G
\]
 for some \(n\geq 0\). Prove that the descending central series \(G=G^0\supseteq G^1\supseteq G^2\supseteq\cdots\) also has length~\(n\), and that \(G^i\subseteq Z_{n-i}(G)\) for \(0\leq i\leq n\).
\def\ProblemHeadingText{Problem 4.}
\item
Let~\(R\) be a ring with~\(1\).
\begin{examsubparts}
\item[\textnormal{(a)}]
Prove that if~\(J\) is a divisible abelian group then \(\operatorname{Hom}_{\mathbb Z}(R,J)\) is a divisible left~\(R\)-module.
\item[\textnormal{(b)}]
Prove that if~\(M\) is a unitary left~\(R\)-module then~\(M\) is isomorphic to a submodule of a divisible left~\(R\)-module.
\end{examsubparts}
\def\ProblemHeadingText{Problem 5.}
\item
Let~\(C\) be a category and let~\(\Lambda\) be a set. For each \(\lambda\in\Lambda\) let \(X_\lambda\) be an object in~\(C\).
\begin{examsubparts}
\item[\textnormal{(a)}]
Give the definition of a coproduct of the collection of objects \(\{X_\lambda\}_{\lambda\in\Lambda}\).
\item[\textnormal{(b)}]
Prove that any two coproducts of the \(\{X_\lambda\}_{\lambda\in\Lambda}\) are isomorphic.
\item[\textnormal{(c)}]
Prove that coproducts always exist in the category of abelian groups.
\end{examsubparts}
\def\ProblemHeadingText{Problem 6.}
\item
Let~\(F\) be a field and let \(F[[X]]\) denote the ring of formal power series over~\(F\).
\begin{examsubparts}
\item[\textnormal{(a)}]
Prove that \(a_0+a_1X+a_2X^2+\cdots\in F[[X]]\) is a unit if and only if \(a_0\neq 0\).
\item[\textnormal{(b)}]
Prove that \(F[[X]]\) is a discrete valuation ring.
\end{examsubparts}
\def\ProblemHeadingText{Problem 7.}
\item
Let~\(R\) be a commutative Noetherian ring with~\(1\) and let \(I\subset R\) be an ideal. Prove that~\(I\) has a primary decomposition.
\def\ProblemHeadingText{Problem 8.}
\item
Let~\(R\) be a left Artinian ring. Prove that the following are equivalent:
\begin{examsubparts}
\item[\textnormal{(a)}]
\(R\) is simple.
\item[\textnormal{(b)}]
\(R\) is primitive.
\item[\textnormal{(c)}]
\(R\cong M_n(D)\) for some \(n\geq 1\) and division ring~\(D\).
\end{examsubparts}
\def\ProblemHeadingText{Problem 9.}
\item
Let~\(R\) be a ring and let~\(M\) be a left~\(R\)-module. Assume that for every nonzero \(x\in M\) we have \(Rx\neq 0\), and that for every submodule \(N\subseteq M\) there is a submodule \(N'\) such that \(M=N\oplus N'\). Prove that~\(M\) is generated by its simple submodules.
\def\ProblemHeadingText{Problem 10.}
\item
This problem has two parts.
\begin{examsubparts}
\item[\textnormal{(a)}]
State the Noether–Skolem theorem.
\item[\textnormal{(b)}]
Let~\(D\) be a division ring such that

\par
Prove that~\(D\) is isomorphic to either \(\mathbb R\), \(\mathbb C\), or \(\mathbb H\).

\par
(Here \(\mathbb R\) is the real numbers, \(\mathbb C\) is the complex numbers, and \(\mathbb H\) is Hamilton’s quaternions.)
\begin{examsubparts}
\item[\textnormal{i.}]
\(\mathbb R\) is a subfield of~\(D\).
\item[\textnormal{ii.}]
Every \(\alpha\in D\) is algebraic over \(\mathbb R\).
\end{examsubparts}
\end{examsubparts}
\def\ProblemHeadingText{Problem 11.}
\item
Let \(m,n\geq 1\). Describe the ring \((\mathbb Z/m\mathbb Z)\otimes_{\mathbb Z}(\mathbb Z/n\mathbb Z)\). In particular, what is the cardinality of this ring?
\end{examproblems}
\end{document}
